%%%%%%%%%%%%%%%%%%%%%%%%%%%%%%%%%%%%%%%%%%%%%%%%%%%%%%%%%%%%%%%%%%%%%%
% LaTeX Template: Curriculum Vitae
%
% Source: http://www.howtotex.com/
% Feel free to distribute this template, but please keep the
% referal to HowToTeX.com.
% Date: July 2011
% 
% Author of Anti-CV: Stephen Fay
% Fork on github https://github.com/dcxSt/anti-cv
% 

\documentclass[paper=a4,fontsize=11pt]{scrartcl} % KOMA-article class
							
\usepackage[english]{babel}
\usepackage[utf8x]{inputenc}
\usepackage[protrusion=true,expansion=true]{microtype}
\usepackage{amsmath,amsfonts,amsthm}     % Math packages
\usepackage{graphicx}                    % Enable pdflatex
\usepackage[svgnames]{xcolor}            % Colors by their 'svgnames'
\usepackage[margin=0.5in]{geometry}
	\textheight=700px                    % Saving trees ;-)
\usepackage{url}

\usepackage{tikz}
\usepackage{kpfonts} % for \varheartsuit , also changes font
\usepackage{color}
\newcommand*\Hs{\ensuremath{\color{red}\varheartsuit}}

%%% Heart-stab command
\newcommand{\heartstab}{%
\begin{tikzpicture}%
\draw [line width=0.20mm, white] (3.5ex,0.4ex) -- (3.5ex,0.4ex);
\Large{$\Hs$};
\draw [line width=0.30mm, black] (-0.7ex,0.7ex) -- (1.2ex,1.2ex);%
\draw [line width=0.50mm, black] (0.8ex,0.65ex) -- (0.55ex,1.45ex);%
\end{tikzpicture}%
}

%%% Squiggly arrow
\newcommand{\squigarr}{
{\Large \color{blue}$\rightsquigarrow$}
}

\usepackage[colorlinks=true,urlcolor=blue]{hyperref}

\frenchspacing              % Better looking spacings after periods
\pagestyle{empty}           % No pagenumbers/headers/footers

%%% Custom sectioning (sectsty package)
%%% ------------------------------------------------------------
\usepackage{sectsty}

\sectionfont{%			            % Change font of \section command
	\usefont{OT1}{phv}{b}{n}%		% bch-b-n: CharterBT-Bold font
	\sectionrule{0pt}{0pt}{-5pt}{3pt}}

%%% Macros
%%% ------------------------------------------------------------
\newlength{\spacebox}
\settowidth{\spacebox}{8888888888}			% Box to align text
\newcommand{\sepspace}{\vspace*{0.5em}}		% Vertical space macro
        
\newcommand{\NameAffiliation}[5]{
  \noindent
  \begin{minipage}[t]{0.45\textwidth}
    \vspace{0pt} % Align at top
    {\fontsize{32}{32} \usefont{OT1}{phv}{m}{n} \textbf{#1}}
  \end{minipage}%
  \hfill
  \begin{minipage}[t]{0.5\textwidth}
    \vspace{0pt} % Align at top
    \raggedleft
    {\large \usefont{OT1}{phv}{m}{n} #2 \\}
    {\large \usefont{OT1}{phv}{m}{n} #3 \\}
    {\large \usefont{OT1}{phv}{m}{n} #4 \\}
    {\large \usefont{OT1}{phv}{m}{n} #5}
  \end{minipage}
}

\newcommand{\NewPart}[1]{\section*{{#1}}}

\newcommand{\SkillsEntry}[2]{      % Same as \PersonalEntry
		\noindent\hangindent=2em\hangafter=0 % Indentation
		\parbox{\spacebox}{        % Box to align text
		\textit{#1}}			   % Entry name (birth, address, etc.)
		\hspace{1.5em} #2 \par}    % Entry value	
		


%%% Begin Document
%%% ------------------------------------------------------------
\begin{document}

%% you can upload a photo and include it here...
%\begin{wrapfigure}{l}{0.5\textwidth}
%	\vspace*{-2em}
%		\includegraphics[width=0.15\textwidth]{picture-of-me1.jpg}
%\end{wrapfigure}

% ALTERNATIVE PICTURE SETTINGS TO THE ABOVE :
%%%%%%%%%%%%%%%%%%%%%%%%%%%%%%%%%%%%%%%%%%%%%%%%
%\iffals
%\begin{figure}
%%  \includegraphics[width=\linewidth]{picture-of-me1.jpg}
%%  \caption{
%    \includegraphics[width=.15\textwidth,left]{picture-of-me1.jpg}
%%    caption{left aligned image}
%%    \label{fig:me}
%\end{figure}
%\fi
%%%%%%%%%%%%%%%%%%%%%%%%%%%%%%%%%%%%%%%%%%%%%%%%

%\MyName{Stephen Fay}

%\MySlogan{Curriculum Vitae}

\NameAffiliation
{Zhuo Liu}
{Department of Earth and Planetary Science, Stanford University}
{450 Jane Stanford Way}
{Bldg. 320, Rm.118}
{Stanford, CA 94305, USA}

\sepspace

%%% Contact Info
%%% ------------------------------------------------------------
\NewPart{Personal Information}{}
\hspace{2em} Email: \url{zliu93@stanford.edu} \hspace{10em} Tel: +1 303 269 1936 

\hspace{1em} Email: \url{liuzhuolz@outlook.com} \hspace{8.95em} Tel: +86 150 8499 8048

\hspace{1em} Linkedin: \url{linkedin.com/in/liuzzzz/} 
\sepspace




%%% Employment
%%% ------------------------------------------------------------
\NewPart{Employment}{}
\hspace{2em} \textbf{Stanford University, Stanford, CA} \hspace{20em} Mar. 2025 – Present

\hspace{1em} Postdoctoral Scholar with Stanford Mineral-X Initiative

\hspace{1em} Principal Investigator: Dr. Jef Caers
\sepspace
%%% ============================================================



%%% Education
%%% ------------------------------------------------------------
\NewPart{Education}{}
\hspace{2em} \textbf{Colorado School of Mines, Golden, CO} \hspace{17.1em} Jan. 2019 – Dec. 2024

\hspace{1em} Department of Geophysics

\hspace{1em} Ph.D., major in Geophysics, advisor: Dr. Yaoguo Li  \hspace{16.5em} GPA: 3.57/4

\hspace{1.2em}Minor in Geology, advisor: Dr. Thomas Monecke

\noindent\hangindent=2.3em\hangafter=0 Dissertation Title: Magnetic inversions constrained with geological information: An investigation using regularized inversion and machine learning approach.

\bigskip
\hspace{1em} \textbf{Colorado School of Mines, Golden, CO} \hspace{16.5em} Sept. 2016 – Dec. 2018

\hspace{1em} Department of Geophysics

\hspace{1em} M.Sc., major in Geophysics, advisor: Dr. Yaoguo Li \hspace{16.5em} GPA: 3.56/4

\noindent\hangindent=2.3em\hangafter=0 Dissertation Title: Effects of anisotropic magnetic susceptibility in data interpretation and its potential in application

\bigskip
\hspace{1em} \textbf{Central South University, Changsha, Hunan, China } \hspace{10.5em} Sept. 2011 – Jun. 2015

\hspace{1em} School of Geoscience and Info-Physics

\hspace{1em} B.Sc., major in Applied Geophysics \hspace{23em} GPA: 3.29/4
\sepspace
%%% ============================================================




%%% Research Interests
%%% ------------------------------------------------------------
\NewPart{Research Interests}{}
\noindent\hangindent=2.3em\hangafter=0 My research focuses mainly on exploiting the information contained in the geophysical data sets and fostering geological interpretation using geophysical data interpretation approaches. 
I am particularly interested in extracting both geophysical and geological information from potential field data through inversion, geology differentiation based on multi-physical properties and anisotropies, and incorporating conceptual prior information into inversion assisted by machine learning approaches. 
Additional research interest is in the topics of turning trained neural networks into inverse operators or assistants of the decision-making process.
\sepspace
%%% ============================================================




%%% Pear-reviewd Paper
%%% ------------------------------------------------------------
\NewPart{Pear-reviewed Paper}{}
\begin{enumerate}
    \item \textbf{Zhuo Liu} and Yaoguo Li, “Connect Geophysical Data Interpretation and Geology Through Inversion for Anisotropic Magnetic Susceptibility.” \textit{Geophysical Prospecting} (2025): e70037.
    \item \textbf{Zhuo Liu}, Yaoguo Li, and Kaijun Xu, “Constrained Simultaneous Recovery of the Depth to Basement and Lateral Susceptibility Variation.” \textit{Geophysical Prospecting} 72, no. 8 (2024): 3008-3025.
    \item \textbf{Zhuo Liu} and Yaoguo Li, “Simultaneous Imaging of Basement Relief and Varying Susceptibility Using Trained Deep Neural Network.” Submitted to \textit{Geophysics}.
\end{enumerate}
\sepspace
%%% ============================================================




%%% Conference Presentation
%%% ------------------------------------------------------------
\NewPart{Conference Presentation}{}
\begin{enumerate}
    \item \textbf{Zhuo Liu*}  and Yaoguo Li. “Simultaneous imaging of basement relief and varying susceptibility in deep-learning approach.” In The International Meeting for Applied Geoscience \& Energy. SEG \& AAPG, 2024.
    \item \textbf{Zhuo Liu*}, Yaoguo Li, and Kaijun Xu. “Magnetotelluric Constrained Simultaneous Recovery of the Depth to Basement and Lateral Susceptibility Variation.” In The International Meeting for Applied Geoscience \& Energy. SEG \& AAPG, 2024
    \item \textbf{Zhuo Liu*} and Yaoguo Li. “Inversion for anisotropy magnetic susceptibility.” In SEG International Exposition and Annual Meeting, p. D031S028R007. SEG, 2020.
    \item \textbf{Zhuo Liu*} and Yaoguo Li. “Effects of anisotropic magnetic susceptibility in data interpretation and potential in application.” In International Workshop on Gravity, Electrical \& Magnetic Methods and Their Applications, Xi'an, China, May 19–22, 2019, pp. 351-354. Society of Exploration Geophysicists and Chinese Geophysical Society, 2019.
    \item \textbf{Zhuo Liu*} and Yaoguo Li. “Effect of anisotropic magnetic susceptibility and potential applications.” In SEG International Exposition and Annual Meeting, pp. SEG-2018. SEG, 2018
\end{enumerate}
\sepspace
%%% ============================================================




%%% Teaching Experience
%%% ------------------------------------------------------------
\NewPart{Teaching Experience}{}
\hspace{1em} \textbf{Colorado School of Mines} \hspace{27.5em} Golden, CO

\hspace{0em} Advanced Gravity and Magnetic Exploration (GPGN 411/511)  \hspace{12.8em} Fall 2023

\hspace{0em} \textbullet \hspace{1em} Teach the inversion section regarding the Tikhonov regularized inversion and Newton method.

\hspace{0em} \textbullet \hspace{1em} Grade homework and exam assignments.

\hspace{0em} Physics of Earth I (GPGN 328)  \hspace{26.9em} Fall 2022

\hspace{0em} \textbullet \hspace{1em} Host review sessions for each section and before exams.

\hspace{0em} \textbullet \hspace{1em} Teach the vector multiplication section.

\hspace{0em} \textbullet \hspace{1em} Grade homework and exam assignments.

\sepspace
%%% ============================================================




%%% Field Experience
%%% ------------------------------------------------------------
\NewPart{Field Experience}{}
\hspace{1em} Economy Geology (GEOL 524) class field trip to Edgar Mine at Idaho Springs, CO, USA \hspace{1.6em} Nov 2019

\hspace{0em} \textbullet \hspace{1em} Map mineral occurrence and structural trend in the tunnel.

\hspace{0em} Hydrothermal Geochemistry (GEOL 513) class field trip to Steamboat Springs, CO, USA \hspace{1.3em} Mar 2019

\hspace{0em} \textbullet \hspace{1em} Sampling geothermal fluid and perform pH, Eh, conductivity, and alkalinity analyses on site.

\hspace{0em} CSEM data collection in Songliao Basin near Biacheng, Jilin, China  \hspace{11.2em} Jul 2017

\hspace{0em} \textbullet \hspace{1em} Collect data and perform data quality examination on site.

\sepspace
%%% ============================================================




%%% SProfessional Skills
%%% ------------------------------------------------------------
\NewPart{Professional Skills}{}
\hspace{1em} \textbf{Software}: Adobe Photoshop, Microsoft Office, Mathematica, Geosoft Oasis Montaj, Esri ArcGIS

\hspace{0em} \textbf{Coding}: MATLAB, Python
\sepspace
%%% ============================================================



    
%%% Honor and Awards
%%% ------------------------------------------------------------
\NewPart{Honor and Awards}{}
\hspace{1em} \textbf{Meng Ersheng Endowed Scholarship } \hspace{23em} Nov 2024

\hspace{0em} By the Department of Geophysics, Colorado School of Mines

\hspace{0em} \textbf{Excellent Graduate of Hunan Province} \hspace{22.7em} Jun 2015

\hspace{0em} By the Education Department of Hunan Province

\hspace{0em} \textbf{Top Ten Student of Central South University} \hspace{19.9em} Jun 2015

\hspace{0em} By Central South University

\hspace{0em} \textbf{Outstanding student} \hspace{30.4em} Nov 2013

\hspace{0em} By the School of Geoscience and Info-Physics

\hspace{0em} \textbf{First-Class Laurel Geophysics Scholarship} \hspace{20.6em} Nov 2013

\hspace{0em} By Laurel Technologies Co Ltd
%%% ============================================================

\end{document}


